\chapter{Heaviest Common Subsequence (HCS)}

\section{Definizione del problema}

\paragraph{Istanza:}
\begin{itemize}
    \item Due sequenze su un alfabeto finito $\Sigma$:
    \[
        X = \langle x_1, x_2, \dots, x_m \rangle, \qquad Y = \langle y_1, y_2, \dots, y_n \rangle
    \]
    \item Una funzione peso definita sui simboli dell'alfabeto:
    \[
        w : \Sigma \to \mathbb{N}^+
    \]
\end{itemize}

\paragraph{Soluzioni ammissibili:}
Tutte le sottosequenze comuni di $X$ e $Y$.

\paragraph{Funzione obiettivo:}
Il peso complessivo $W(Z)$ di una sottosequenza comune $Z = \langle z_1, z_2, \dots, z_k \rangle$, definito come la somma dei pesi dei suoi elementi:
\[
    W(Z) = \sum_{l=1}^k w(z_l)
\]

\paragraph{Soluzione ottima:}
Una sottosequenza comune $Z = \text{HCS}(X, Y)$ che massimizza il peso totale $W(Z)$. La dimensione del problema principale è $(m, n)$.

\section{Definizione del generico sottoproblema}

Per ogni coppia di indici $i \in \{0, 1, \dots, m\}$ e $j \in \{0, 1, \dots, n\}$, consideriamo i prefissi:
\[
    X_i = \langle x_1, \dots, x_i \rangle \quad (X_0 = \langle \rangle), \qquad Y_j = \langle y_1, \dots, y_j \rangle \quad (Y_0 = \langle \rangle)
\]

\paragraph{Sottoproblema di dimensione $(i, j)$:} Trovare una più pesante sottosequenza comune tra i prefissi $X_i$ e $Y_j$.
\begin{itemize}
    \item \textbf{Numero totale di sottoproblemi:} $(m+1)(n+1)$.
    \item \textbf{Coefficiente $c_{i,j}$ (valore ottimo):} Peso totale della più pesante sottosequenza comune tra $X_i$ e $Y_j$.
    \item \textbf{Valore ottimo del problema principale:} $c_{m,n}$.
\end{itemize}

\section{Equazioni di ricorrenza}

\subsection{Casi base}
Se almeno una delle due sequenze è vuota ($i = 0 \lor j = 0$), l'unica sottosequenza comune è la sequenza vuota $\langle \rangle$, che ha peso nullo:
\begin{equation}
    c_{i,j} = 0 \qquad \text{se } i = 0 \lor j = 0
\end{equation}

\subsection{Passo ricorsivo}
Per $i > 0$ e $j > 0$, l'equazione si articola in base al confronto tra gli ultimi caratteri dei prefissi $x_i$ e $y_j$:
\begin{enumerate}
    \item \textbf{Se $x_i \ne y_j$:} I caratteri non possono appartenere entrambi a un match comune finale. Il valore ottimo corrisponde al massimo tra le soluzioni dei due sottoproblemi ridotti:
    \begin{equation}
        c_{i,j} = \max\{c_{i-1,j}, \, c_{i,j-1}\}
    \end{equation}
    \item \textbf{Se $x_i = y_j$:} Il carattere comune entra a far parte della sottosequenza e contribuisce con il proprio peso $w(x_i)$:
    \begin{equation}
        c_{i,j} = c_{i-1,j-1} + w(x_i)
    \end{equation}
\end{enumerate}

\section{Algoritmo iterativo Bottom-Up}

L'algoritmo calcola sistematicamente i valori della matrice $c[0 \dots m, \, 0 \dots n]$ procedendo per righe o per colonne a partire dai casi base.

\begin{algorithm}[H]
\caption{Calcolo della matrice HCS}\label{alg:hcs}
\begin{algorithmic}[1]
\Procedure{HCS}{$X, Y, w$}
    \State $m \gets |X|$
    \State $n \gets |Y|$
    \State $c[0 \dots m, \, 0 \dots n] \gets \text{nuova matrice}$
    \Statex \Comment{Inizializzazione dei casi base}
    \For{$i \gets 0 \textbf{ to } m$}
        \State $c[i, 0] \gets 0$
    \EndFor
    \For{$j \gets 0 \textbf{ to } n$}
        \State $c[0, j] \gets 0$
    \EndFor
    \Statex \Comment{Passo ricorsivo}
    \For{$i \gets 1 \textbf{ to } m$}
        \For{$j \gets 1 \textbf{ to } n$}
            \If{$x_i \ne y_j$}
                \State $c[i, j] \gets \max(c[i-1, j], \, c[i, j-1])$
            \Else
                \State $c[i, j] \gets c[i-1, j-1] + w(x_i)$
            \EndIf
        \EndFor
    \EndFor
    \State \Return $c$
\EndProcedure
\end{algorithmic}
\end{algorithm}

\paragraph{Analisi della complessità di \textsc{HCS}:}
\begin{itemize}
    \item \textbf{Tempo:} $\Theta(m \cdot n)$, determinato dai due cicli for annidati in cui ogni cella viene calcolata in tempo costante $\Theta(1)$.
    \item \textbf{Spazio:} $\Theta(m \cdot n)$ per memorizzare la tabella dei coefficienti $c$.
\end{itemize}

\newpage
\section{Algoritmo di ricostruzione}

La procedura \textsc{HCSR} risale la matrice a ritroso partendo da $(m, n)$ fino a raggiungere un caso base ($i=0 \lor j=0$). La stampa in post-ordine (eseguita dopo la chiamata ricorsiva) garantisce che i caratteri della soluzione ottima vengano stampati nell'ordine corretto da sinistra verso destra.

\begin{algorithm}[H]
\caption{Ricostruzione della soluzione HCS}\label{alg:hcsr}
\begin{algorithmic}[1]
\Procedure{HCSR}{$c, i, j, X, Y$}
    \If{$i == 0 \lor j == 0$}
        \State \Return
    \EndIf
    \If{$x_i \ne y_j$}
        \If{$c[i-1, j] \ge c[i, j-1]$}
            \State \Call{HCSR}{$c, i-1, j, X, Y$}
        \Else
            \State \Call{HCSR}{$c, i, j-1, X, Y$}
        \EndIf
    \Else
        \State \Call{HCSR}{$c, i-1, j-1, X, Y$}
        \State $\textbf{print}(x_i)$
    \EndIf
\EndProcedure
\end{algorithmic}
\end{algorithm}

\paragraph{Analisi della complessità di \textsc{HCSR}:}
\begin{itemize}
    \item \textbf{Tempo:} $O(m + n)$, poiché a ogni passo ricorsivo viene decrementato almeno uno dei due indici $i$ o $j$.
    \item \textbf{Spazio ausiliario:} $O(m + n)$ dovuto alla profondità massima dello stack delle chiamate ricorsive.
\end{itemize}